% pi_prose_1125.tex -- PI text of 2026-11-25 for paper/journal/v2
%
% HOW TO USE (Claude Code)
%   \bn{...}  a number that must come from the build. The value shown is the one in checkpoint 98b812b;
%             replace each with its macro. If the final read-out moves it, the macro wins.
%   \ct{...}  a citation. Map each name to its bib key. Cite only after reading the paper.
%   \chk{...} a clause that depends on a row or on a paper I could not verify. Resolve each one;
%             if the row fails or the paper does not say it, do not reword: send me the sentence.
%   Abstract and Introduction are installed as written. Discussion is a draft to complete after part 1.
%   The two model paragraphs show the voice for Methods and Results. Imitate the voice, not the wording.
%
% This file compiles on its own as a preview. In v2, copy the text between the markers.

\documentclass[11pt,a4paper]{article}
\usepackage[margin=3.2cm]{geometry}
\usepackage[T1]{fontenc}
\usepackage{lmodern}
\usepackage{xcolor}
\usepackage{microtype}
\newcommand{\bn}[1]{\textcolor{blue!70!black}{#1}}
\newcommand{\ct}[1]{\textcolor{black!55}{[#1]}}
\newcommand{\chk}[1]{\textcolor{red!70!black}{#1}}
\newcommand{\bb}{b\bar{b}}
\newcommand{\cc}{c\bar{c}}
\title{\Large Label granularity in supervised pretraining of jet foundation models\\[4pt]
\normalsize Prose preview, 25 November 2026: abstract, introduction, discussion, two model paragraphs}
\author{}
\date{}
\begin{document}
\maketitle
\noindent{\small Blue: a number the build must supply. Grey: a citation to map. Red: a clause to verify before it is printed.}

%%%%%%%%%%%%%%%%%%%%%%%%%%%%%%%%%%%%%%%%%%%%%%%%%%%%%%%%%%%%%%%%%%%%%%%%%%%%%%%%%%%%%%%%%%%%%%%%%%%
% BEGIN ABSTRACT
\begin{abstract}
\noindent
When a foundation model is pretrained by supervised classification, its classes determine which
distinctions the network must learn, yet how fine they need to be remains an open question.
Jets, the collimated sprays of particles produced in high-energy collisions, offer a clean test,
because each simulated jet carries an exact label that can be coarsened without changing the
training data.
We pretrain Sophon, a foundation model for jets, on the same simulated jets under four nested label sets,
from 188 classes down to 17, which successively merge decays differing only in the flavor of the
particles produced.
As the labels coarsen, the frozen representation retains less of what they no longer distinguish:
a linear classifier trained on it identifies the 188 classes with \bn{43\%} accuracy for the
finest labels but \bn{33\%} for the coarsest.
Random groupings into 17 classes, however, still reach \bn{42\%}, so the loss depends on which
classes are merged and not only on how many.
Fine-tuning narrows the gap to about \bn{2} percentage points without closing it, and the
coarsest labels also remain behind on a second multi-class task, while two binary tagging
benchmarks show no significant difference.
Beyond the labels, an auxiliary regression of the jet mass sharpens the mass resolution of the
frozen representation at no significant cost in its accuracy.
Fine labels therefore matter most for tasks that rely on the distinctions they preserve, and
when the model is used without fine-tuning.
\end{abstract}
% END ABSTRACT

%%%%%%%%%%%%%%%%%%%%%%%%%%%%%%%%%%%%%%%%%%%%%%%%%%%%%%%%%%%%%%%%%%%%%%%%%%%%%%%%%%%%%%%%%%%%%%%%%%%
% BEGIN INTRODUCTION
\section{Introduction}

Foundation models are pretrained once on broad data and then adapted to many tasks
\ct{foundation-model reference}.
When the pretraining is supervised, the label set is itself a design choice: the same examples can
be divided into few classes or many, and that division determines which distinctions the network
is rewarded for learning.
In image classification, relabeling a fixed set of images along a class hierarchy has shown that
coarser labels give features that transfer less well, and more so when the features are used
frozen than after fine-tuning \ct{Huh et al.; Hong et al.}.
Such hierarchies, however, are human taxonomies, and it is seldom clear what information a
given merge removes from the labels.

Particle physics allows the question to be put more sharply, because there the labels come from
simulation.
At the Large Hadron Collider, quarks, gluons and the decay products of heavy particles reach the
detector as jets, collimated sprays of tens of particles, and identifying what produced a jet,
known as jet tagging, is one of the most developed uses of machine learning in the field
\ct{2--6}.
Each simulated jet comes with an exact record of the process that produced it, so its label can
be made as fine as the simulated physics allows.
Sophon \ct{Sophon}, a foundation model for jets \ct{12--16} built on the Particle Transformer
\ct{ParT}, uses this freedom: it is pretrained to sort the jets of the JetClass-II dataset
\ct{JetClass-II} into 188 classes, which specify the particle that decayed and what it decayed
into.

Whether such fine labels make the pretrained model better has not been tested in a controlled way.
In the comparisons published so far the training jets change together with the labels, so that a
model trained on fewer classes has also seen fewer kinds of jets
\ct{Sophon; Li thesis; Uslu et al.; SuperJEPA}, and the two effects cannot be told apart.

Here we change the labels and nothing else.
Since every JetClass-II jet records its full decay, the same jets can be relabeled at any
granularity, and we pretrain the Sophon network on them with four nested label sets of 188, 162,
43 and 17 classes (Table~1).
Each coarser set merges classes that differ only in flavor, the type of quark or lepton produced:
in the coarsest, for example, a decay to two bottom quarks and a decay to two charm quarks both
become a decay to two quarks or gluons.
Random groupings of the classes into 17 then separate the effect of which classes are merged from
that of how many remain, and an auxiliary regression of the jet mass adds a target that the labels
do not contain.

We evaluate each pretrained model in the two ways a foundation model is used.
With the network frozen, a linear classifier trained on its features (a linear probe) measures
what the representation holds, both for the 188 pretraining classes and for pairs of classes that
differ in a single flavor.
With the network fine-tuned, we measure what survives adaptation: on the pretraining classes, on
the ten classes of JetClass \ct{ParT}, which come from a different simulation, and on two public
binary benchmarks, the identification of jets from top quarks (top tagging) \ct{top landscape}
and the discrimination of quark jets from gluon jets \ct{26--28}.
We also repeat the benchmark with which Sophon was introduced, the tagging of a hypothetical
heavy particle $X$ that decays to $\bb$ or to $bs$
\chk{[one clause: what Sophon does for each decay and what we do]}.
Repeated pretraining runs give every comparison an uncertainty.

We find that the frozen representation loses part of what the labels no longer distinguish.
A linear probe identifies the 188 classes with \bn{43\%} accuracy on the features of the 188-class
model and \bn{33\%} on those of the 17-class model, pairs of classes that a label set merges
are separated less well, and \chk{the resonance benchmark loses signal efficiency for both
decays}.
The number of classes does not by itself set this loss, since random groupings into 17 classes
still reach \bn{42\%}.
Fine-tuning on $10^6$ jets recovers part of the loss: the 17-class model remains \bn{2.3}
percentage points behind the 188-class model on the pretraining classes and \bn{0.5} points behind
on JetClass, the 43-class model is also significantly behind, and on the two binary benchmarks no
difference between the label sets is significant
\chk{[scope to background rejection if another metric passes at the full sets]}.
The mass regression, finally, sharpens the mass resolution of the frozen features at no
significant cost in their accuracy.
Fine labels thus matter most for tasks that rely on the distinctions they preserve, and when the
pretrained model is used without fine-tuning.
% END INTRODUCTION

%%%%%%%%%%%%%%%%%%%%%%%%%%%%%%%%%%%%%%%%%%%%%%%%%%%%%%%%%%%%%%%%%%%%%%%%%%%%%%%%%%%%%%%%%%%%%%%%%%%
% BEGIN DISCUSSION (draft: complete after part 1; every clause keeps its row)
\section{Discussion and conclusion}

The results describe a representation shaped by what its labels distinguish, and not simply by
how many labels there are.
Pairs of classes that a label set merges lose separation in the frozen features, and at 43 classes
some pairs that the labels still keep apart lose it as well.
Random groupings of the same size as the coarsest set nevertheless retain most of the accuracy
on the pretraining classes, so the number of classes alone does not determine the loss.
That result should not be read as a recipe for label design, since on JetClass the random
groupings offer no advantage over the 17-class set
\chk{and on the two binary benchmarks their frozen features are the weakest of any pretrained
model [state only what the bootstrap rows support]}.
For an analysis that reads a discriminant from frozen features, or sums class scores into one as
in the resonance benchmark, the practical consequence is that the classes it must tell apart
should be distinct in the pretraining labels.

These findings agree with relabeling experiments in image classification and extend them.
There too, coarser labels cost more for frozen features than after fine-tuning \ct{Hong et al.},
\chk{and classes merged along a semantic hierarchy remain more entangled than classes merged at
random} \ct{\chk{Yang, Steinhardt and Hu}}.
What simulation adds is control: the labels are exact, the training examples and their order are identical for every label set,
and each merge removes a distinction that can be named, so the loss can be traced to individual
pairs of classes instead of being inferred from an average over a taxonomy.

Fine-tuning reduces the loss without removing it: with $10^6$ labeled jets the 17-class model
stays \bn{2.3} points behind the 188-class model on
the pretraining classes and \bn{0.5} points behind on JetClass, against gains from pretraining of
\bn{6.0} and \bn{2.1} points, so a coarse label set forfeits an appreciable share of what pretraining
provides even when the whole network is adapted.
On top tagging and quark--gluon discrimination no difference between the label sets is
significant, with a run-to-run spread of \chk{[size, from the build]} in background rejection.
The mass regression acts on a different axis: it supplies a quantity that no label set contains
and improves the mass resolution of the frozen features at no significant cost in their accuracy,
although after fine-tuning the models pretrained with it are slightly less accurate, by
\chk{[size at 162 and at 17 classes]}.

These conclusions rest on one architecture, with 2.3~million parameters, and on one simulated
dataset.
Each step of the hierarchy also removes several distinctions at once, so the cost of a single
distinction is established only through pairs of classes, and on the two binary benchmarks several
inputs used in pretraining are absent, which limits what the frozen features can show there.
Within these limits, a label set built for transfer should keep apart the classes that later
analyses will need to separate, and this matters most when the pretrained model is used without
fine-tuning; a quantity that no labels carry, such as the jet mass, can be added as an auxiliary
target.
% END DISCUSSION

%%%%%%%%%%%%%%%%%%%%%%%%%%%%%%%%%%%%%%%%%%%%%%%%%%%%%%%%%%%%%%%%%%%%%%%%%%%%%%%%%%%%%%%%%%%%%%%%%%%
% MODEL PARAGRAPHS: the voice for Methods and Results. Facts and numbers to be checked by the writer.
\section*{Model paragraph for Methods: background for a reader outside jet physics}

A jet is the detector signature of an energetic quark or gluon, or of a heavy particle whose
decay products are collimated by its momentum.
For a machine-learning model it is an unordered set of a few tens of reconstructed particles, each
described here by 17 numbers: its momentum relative to the jet, its electric charge and type and,
for charged particles, the displacement of its track from the collision point, which reveals the
decays of the long-lived hadrons that contain bottom or charm quarks.
The classes of JetClass-II name what the jet came from: either a heavy resonance, labeled by the
quarks, gluons and leptons into which it decayed, or the ordinary production of quarks and gluons
by the strong interaction (QCD jets), which is the background in most analyses.
The two binary benchmarks are standard in the field: top tagging asks whether a jet contains the
decay of a top quark, and quark--gluon discrimination whether it was initiated by a quark or by a
gluon.

\section*{Model paragraph for Results: the pretraining classes}

Coarser labels leave the frozen features less able to tell the pretraining classes apart (Table~2).
A linear probe on the features of the 188-class model identifies the 188 classes with an accuracy
of \bn{43.5\%}, and the 162-class model, which differs only in treating all QCD jets as a single
class, is \bn{0.2} points lower \chk{[a loss confined to the QCD classes it merges: check its row]}.
The accuracy then falls to \bn{36.1\%} for the 43-class model and to \bn{33.4\%} for the 17-class
model, the two sets that merge flavors within the resonance decays.
These values should be read against two references: the same probe reaches \bn{10.2\%} on an
untrained network, and the model released with Sophon gives \bn{43.4\%}, in agreement with our
retraining of it \chk{[add the majority-class and pretrained-output-layer accuracies from part 1]}.

\end{document}
